\section{M\"{o}bius Group}
\label{sec:mobius-group}

This section is about a fantastically interesting group. It is \emph{almost} a group of bijections on \(\CC\), but we need to formally define a `point at infinity'.

\subsection{M\"{o}bius Transformations}

\begin{definition}[Riemann Sphere]
    \label{def:riemann-sphere}%
    The \emph{Riemann Sphere} is
    \[
        \CC\compactify = \CC \cup \brce{\infty}.
    \]
\end{definition}

The geometric explanation of a `sphere' is as follows. Consider the map defined in \autoref{fig:stereographic-projection-riemann-sphere}, the \emph{stereographic projection}
\[
    \functiondomain{\pi}{\CC}{S^2}.
\]

\begin{figure}[ht!]
    \centering
    \caption{Stereographic Projection for the Riemann Sphere}
    \label{fig:stereographic-projection-riemann-sphere}
\end{figure}

This identifies \(\CC\) with all the points of the sphere, except for the `North Pole' of the sphere, which corresponds to \(\infty\).

\begin{remark}
    In geometry terms, we may also think of this as the \emph{complex projective line}, \(\CP^{1}\).
\end{remark}

\begin{definition}[M\"{o}bius Transformation]
    \label{def:mobius-transformation}%
    Let \(a, b, c, d \in \CC\), and \(ad - bc \neq 0\). We define the corresponding \emph{M\"{o}bius Transformation} \(\functiondomain{\mu}{\CC\compactify}{\CC\compactify}\) as follows.

    \begin{itemize}
        \item If \(c \neq 0\), then
        \[
            \functionmap{\mu}{z}{\begin{cases}
                \frac{az + b}{cz + d}, & z \in \CC \setminus \brce{-\frac{d}{c}}, \\
                \infty, & z = - \frac{d}{c} \\
                \frac{a}{c}, & z = \infty.
            \end{cases}}
        \]
        \item If \(c = 0 \), then
        \[
            \functionmap{\mu}{z}{\begin{cases}
                \frac{az + b}{d}, & z \in \CC,\\
                \infty, & z = \infty.
            \end{cases}}
        \]
    \end{itemize}

    We usually just write
    \[
        \mu\para{z} = \frac{az + b}{cz + d}
    \]
    and interpret the different cases appropriately.

    The set of all M\"{o}bius Transformations is denoted as \(\mobius\).
\end{definition}

We compute the composition of two M\"{o}bius transformations. Let
\[
    \mu_{1} \para{z} = \frac{a_1 z + b_1}{c_1 z + d_1}, \mu_{2} \para{z} = \frac{a_2 z + b_2}{c_2 z + d_2},
\]
then
\[
    \mu_{1} \circ \mu_{2} \para{z} = \frac{\para{a_1 a_2 + b_1 c_2} z + \para{a_1 b_2 + b_1 d_2}}{\para{c_1 a_2 + d_1 c_2} z + \para{c_1 b_2 + d_1 d_2}}.
\]

We may compare this with matrix multiplication:
\[
    \begin{pmatrix}
        a_1 & b_1 \\ c_1 & d_1
    \end{pmatrix}
    \begin{pmatrix}
        a_2 & b_2 \\ c_2 & d_2
    \end{pmatrix}
    =
    \begin{pmatrix}
        a_1 a_2 + b_1 c_2 & a_1 b_2 + b_1 d_2 \\
        c_1 a_2 + d_1 c_2 & c_1 b_2 + d_1 d_2
    \end{pmatrix},
\]
and using the fact that the determinant is multiplicative, we see that composition does define a binary operation on \(\mobius\).

\begin{theorem}[M\"{o}bius Group]
    \label{thm:mobius-group}%
    \(\para{\mobius, \circ, \identity_{\CC\compactify}}\) is a group.
\end{theorem}

\begin{proof}
    Composition of functions is associative. The identity
    \[
        \functionmap{\identity_{\CC\compactify}}{z}{z = \frac{1 \cdot z + 0}{0 \cdot z + 1}}
    \]
    is the identity.

    For \(\functionmap{\mu}{z}{\frac{az + b}{cz + d}}\), let \(\functionmap{\nu}{z}{\frac{dz - b}{-cz + a}}\), then
    \[
        \mu \nu = \frac{\para{ad - bc}{z}}{ad - bc} = z
    \]
    gives the identity. (Compare this to inverse of matrices.)
\end{proof}

There is a na\"{i}ve group action from those mappings: \(\mobius \actson \CC\compactify\) by M\"{o}bius transformations.

\begin{remark}
    The previous observation with matrix multiplication shows that the na\"{i}ve mapping is a surjective homomorphism \(\functiondomain{\phi}{\GL_{2} \para{\CC}}{\mobius}\). In fact, we see that the restriction \(\functiondomain{\rightbar{\phi}_{\SL_{2} \para{\CC}}}{\SL_{2} \para{\CC}}{\mobius}\) is surjective as well, and the kernels are
    \[
        \ker \phi = \set{\lambda \identityM}{\lambda \in \CC},
    \]
    and
    \[
        \ker \rightbar{\phi}_{\SL_{2} \para{\CC}} = \brce{\pm \identityM}.
    \]

    In the future, we see by taking quotients, we have
    \[
        \mobius \isomorphic \GL_{2}\para{\CC} / \set{\lambda \identityM}{\lambda \in \CC} = \PGL_{2} \para{\CC},
    \]
    and
    \[
        \mobius \isomorphic \SL_{2}\para{\CC} / \brce{\pm \identityM} = \PSL_{2} \para{\CC}.
    \]

    This shows \(\PGL_{2} \para{\CC} \isomorphic \PSL_{2} \para{\CC}\), and provides actions \(\PGL_{2} \para{\CC} \actson \CP^{1}\), \(\PSL_{2} \para{\CC} \actson \CP^{1}\).
\end{remark}

One important way to study M\"{o}bius transformations is via their \emph{fixed points}.

\begin{definition}[Fixed Point]
    \label{def:fixed-point}%
    Suppose \(\functiondomain{f}{X}{X}\) is a permutation. Then any \(x \in X\) such that \(f\para{x} = x\) is called a \emph{fixed point} of \(f\). The set of all fixed points is denoted as
    \[
        \Fix\para{f} = \set{x \in X}{f\para{x} = x}.
    \]
\end{definition}

Recall that for \(\Isom\para{\CC}\), there is \autoref{lem:three-point-lemma} (Three--Point Lemma). There is a similar lemma for the M\"{o}bius group.

\begin{lemma}[Three--Point Lemma for \(\mobius\)]
    \label{lem:three-point-mobius}%
    If \(\mu \in \mobius\) fixes three distinct points \(w_1, w_2, w_3 \in \CC\compactify\), then \(\mu = \identity\).
\end{lemma}

\begin{proof}
    Let \(\mu\para{z} = \frac{az + b}{cz + d}\). Then a fixed point \(w_i\) satisfies
    \[
        w_i = \frac{a w_i + b}{c w_i + d}.
    \]

    We consider the following cases.

    \begin{itemize}
        \item If there is some \(w_i = \infty\), say \(w_1\), then \(c = 0\). Hence, \(w_2\) and \(w_3\) are such that
        \[
            w_i = \frac{a w_i + b}{d}
        \]
        and hence
        \[
            \para{a - d} w_i + b = 0.
        \]

        This gives a linear equation with \(\geq 2\) roots, so must be trivial. So \(a = d, b = 0\), and \(\mu\para{z} = z\) as desired.

        \item If all \(w_i\)s are not \(\infty\), then \(w_1, w_2\) and \(w_3\) satisfy
        \[
            w_i = \frac{a w_i + b}{c w_i + d}
        \]
        and this rearranges to
        \[
            c w_i^2 + \para{d - a} w_i - b = 0.
        \]

        This is a quadratic with \(\geq 3\) roots, so must be trivial, and hence \(b = c = 0\) and \(a = d\), so \(\mu\para{z} = z\) as desired.
    \end{itemize}
\end{proof}

We can see that any M\"{o}bius transformation has at least \(1\) fixed point in \(\CC\compactify\). In fact, a M\"{o}bius transformation can have \(1, 2\) or \(\infty\)-many fixed points (in the case where it has at least \(3\), it must be the identity and hence trivial).

\begin{examples*}[M\"{o}bius Transformations and Fixed Points]
    \begin{enumerate}
        \item The map
        \[
            \mu\para{z} = z + 1 = \frac{1 \cdot z + 1}{0 \cdot z + 1}
        \]
        fixes \(\infty\):
        \[
            \Fix\para{\mu} = \brce{\infty}.
        \]

        \item If
        \[
        \mu\para{z} = 2z = \frac{2 \cdot z + 0}{0 \cdot z + 1},
        \]
        then \(\mu\) fixes \(\brce{0, \infty}\):
        \[
            \Fix\para{\mu} = \brce{0, \infty}.
        \]
    \end{enumerate}
\end{examples*}

\begin{lemma}[Triple Transitivity]
    \label{lem:triple-transitive}%
    For any triples of distinct points \(z_1, z_2, z_3 \in \CC\compactify\) and \(w_1, w_2, w_3 \in \CC\compactify\), there is a \(\mu \in \mobius\) such that
    \[
        \mu\para{z_i} = w_i
    \]
    for \(i = 1, 2, 3\).
\end{lemma}

\begin{proof}
    Let
    \[
        \functionmap{\alpha \in \mobius}{z}{\frac{z - z_1}{z - z_3} \cdot \frac{z_2 - z_3}{z_2 - z_1},}
    \]
    and we observe that
    \begin{align*}
        z_1 & \mapsto 0, \\
        \alpha: z_2 & \mapsto 1, \\
        z_3 & \mapsto \infty.
    \end{align*}

    Similarly, we can write down \(\beta \in \mobius\) such that
    \begin{align*}
        w_1 & \mapsto 0, \\
        \beta: w_2 & \mapsto 1, \\
        w_3 & \mapsto \infty.
    \end{align*}

    Then the map \(\mu = \beta\inverse \circ \alpha \in \mobius\) gives a construction as desired.
\end{proof}

From \autoref{lem:three-point-mobius} (Three--Point Lemma for the M\"{o}bius Group), we see that the \(\mu\) given by \autoref{lem:triple-transitive} is \emph{unique}. We say the action \(\mobius \actson \CC\compactify\) is \textbf{sharply triply transitive}.

This property allows us to give the following definition, which will show significance later when the geometry of the Riemann Sphere is discussed.

\begin{definition}[Cross Ratio]
    \label{def:cross-ratio}%
    Let \(z_1, z_2, z_3, z_4 \in \CC_{\infty}\) be distinct. Since \(\mobius \actson \CC\compactify\) sharply triply transitively, let \(\alpha \in \mobius\) be the unique M\"{o}bius transformation such that
    \[
        \alpha\para{z_1} = 0, \alpha\para{z_2} = 1, \alpha\para{z_3} = \infty,
    \]
    and the \emph{cross ration} of \(z_1, z_2, z_3, z_4\) is then given by
    \[
        \brac{z_1, z_2, z_3, z_4} = \alpha\para{z_1} = \frac{z_4 - z_1}{z_4 - z_3} \cdot \frac{z_2 - z_3}{z_2 - z_1}.
    \]
\end{definition}

Similar to what we did in dihedral groups, we may use \autoref{lem:three-point-mobius} (Three--Point Lemma for the M\"{o}bius Group) to find a generating set for \(\mobius\).

\begin{proposition}[Generating Set for \(\mobius\)]
    \label{prop:generating-set-mobius}%
    \(\mobius\) is generated by the set of elements of the following \(3\) forms:
    \begin{enumerate}
        \item \(\functionmap{\alpha_{a}}{z}{az}\) for \(a \neq 0, a \in \CC\);
        \item \(\functionmap{\beta_{b}}{z}{z + b}\) for \(b \in \CC\);
        \item \(\functionmap{\gamma}{z}{\frac{1}{z}}\).
    \end{enumerate}
\end{proposition}

Instead of showing this by explicit construction by elementary algebra, we show this similarly to what we did for dihedral groups.

\begin{proof}
    Let \(\mu \in \mobius\), and \(z_1 = \mu\para{0}\), \(z_2 = \mu\para{1}\), and \(z_3 = \mu\para{\infty}\). We `decompose' \(\mu\) by the following three steps.

    \begin{enumerate}
        \item Construct \(\mu_1\) such that \(\mu_1 \para{z_3} = \infty\). If \(z_3 = \infty\), we take \(\mu_1 = \identity_{\CC\compactify}\), or if \(z_3 \neq \infty\), take
        \[
            \mu_{1}\para{z} = \frac{1}{z - z_3} = \gamma \circ \beta_{- z_3} \para{z}.
        \]

        Now, let \(z_1' = \mu_1 \para{z_1}\) and \(z_2' = \mu_1 \para{z_2}\).

        \item Let \(\mu_2 = \beta_{- z_1'}\), i.e.
        \[
            \mu_2\para{z} = z - z_1'.
        \]

        Then \(\mu_2\para{\infty} = \infty\) and \(\mu_2\para{z_1'} = 0\).

        By the construction so far,
        \[
            \mu_2 \circ \mu_1 \para{z_3} = \infty, \mu_2 \circ \mu_1 \para{z_1} = 0.
        \]

        Let \(z_2'' = \mu_2 \circ \mu_1 \para{z_2} \neq 0\).

        \item Let \(\mu_3 = \alpha_{\frac{1}{z_2''}}\). By construction,
        \begin{align*}
            \mu_3 \circ \mu_2 \circ \mu_1 \para{z_1} &= 0, \\
            \mu_3 \circ \mu_2 \circ \mu_1 \para{z_2} &= 1, \\
            \mu_3 \circ \mu_2 \circ \mu_1 \para{z_3} &= \infty,
        \end{align*}
        so \(\para{\mu_3 \circ \mu_2 \circ \mu_1}\inverse = \mu_1\inverse \circ \mu_2\inverse \circ \mu_3\inverse\) sends
        \begin{align*}
            0 &\mapsto z_1, \\
            1 &\mapsto z_2, \\
            \infty &\mapsto z_3,
        \end{align*}
        which are the same as \(\mu\). Therefore, by \autoref{lem:three-point-mobius} (Three--Point Lemma for the M\"{o}bius Group), we have
        \[
            \mu = \mu_1\inverse \circ \mu_2\inverse \circ \mu_3\inverse,
        \]
        as desired.
    \end{enumerate}
\end{proof}

\subsection{Lines and Circles}
\label{sec:lines-and-circles}

We consider the geometry that M\"{o}bius maps preserve.

\begin{definition}[Circle in \(\CC\compactify\)]
    \label{def:circle-riemann}%
    A \emph{circle} in \(\CC\compactify\) is:
    \begin{enumerate}
        \item a Euclidean circle in \(\CC\), defined as
        \[
            \set{z \in \CC}{\abs{z - c} = r}
        \]
        for \(c \in \CC, r > 0\); or
        \item \(L \cup \brce{\infty}\) where \(L\) is a Euclidean straight line in \(\CC\), defined as
        \[
            \set{z \in \CC}{\abs{z - a} = \abs{z - b}}
        \]
        for \(a, b \in \CC\).
    \end{enumerate}
\end{definition}

\begin{figure}[ht!]
    \centering
    \caption{Circles in \(\CC\compactify\)}
    \label{fig:circles-in-riemann-sphere}
\end{figure}

To visualise this, a circle in \(\CC\compactify\) gives a Euclidean straight line in \(\CC\) when it is a circle passing through the `North Pole' of the Riemann sphere.

\begin{theorem}[M\"{o}bius Transformations Preserve Circles]
    \label{thm:mobius-preserve-circles}%
    If \(C \subseteq \CC\compactify\) is a circle and \(\mu \in \mobius\), then \(\mu\para{C}\) is also a circle.
\end{theorem}

\begin{proof}
    Recall that \(\mobius\) is generated by
    \begin{enumerate}
        \item \(\functionmap{\alpha_{a}}{z}{az}\) for \(a \neq 0, a \in \CC\);
        \item \(\functionmap{\beta_{b}}{z}{z + b}\) for \(b \in \CC\);
        \item \(\functionmap{\gamma}{z}{\frac{1}{z}}\).
    \end{enumerate}

    This is easy for \(\alpha_a\) and \(\beta_b\), so we verify for \(\gamma\).
    \begin{itemize}
        \item If \(C\) is a Euclidean circle, \(\gamma\para{C}\) has equation
        \begin{align*}
            \abs{\frac{1}{z} - c} = r &\iff \abs{\frac{1}{z} - c}^2 = r^2 \\
            &\iff \para{\frac{1}{z} - c} \para{\frac{1}{\conjugate{z}} - \conjugate{c}} = r^2 \\
            &\iff \frac{1}{\abs{z}^2} - \frac{c}{\conjugate{z}} - \frac{z}{\conjugate{c}} + \abs{c}^2 - r^2 = 0 \\
            &\iff \para{\abs{c}^2 - r^2} \abs{z}^2 - cz - \conjugate{c} \conjugate{z} + 1 = 0.
        \end{align*}
        \begin{itemize}
            \item If \(\abs{c}^2 = r^2\), i.e. when \(C\) passes through the origin, this is equivalent to
            \begin{align*}
                \abs{\frac{1}{z} - c} = r &\iff cz + \conjugate{c} \conjugate{z} = 1 \\
                &\iff \frac{z}{\conjugate{c}} + \frac{\conjugate{z}}{c} = \frac{1}{\abs{c}^2} \\
                &\iff \abs{z}^2 = \abs{z}^2 - \frac{z}{\conjugate{c}} - \frac{\conjugate{z}}{c} + \frac{1}{\abs{c}^2} \\
                &\iff \abs{z} = \abs{z - \frac{1}{c}}
            \end{align*}
            gives a straight line as required.
            \item Otherwise, \(\abs{c}^2 \neq r^2\), and the equation becomes
            \begin{align*}
                \abs{\frac{1}{z} - c} = r &\iff \abs{z}^2 - \para{\frac{c}{\abs{c}^2 - r^2}} z - \para{\frac{\conjugate{c}}{\abs{c}^2 - r^2}} \conjugate{z} + \frac{1}{\abs{c}^2 - r^2} = 0 \\
                &\iff \abs{z - \frac{\conjugate{c}}{\abs{c}^2 - r^2}}^2 = \frac{\abs{c}^2}{\para{\abs{c}^2 - r^2}^2} - \frac{1}{\abs{c}^2 - r^2} \\
                &\iff \abs{z - \frac{\conjugate{c}}{\abs{c}^2 - r^2}}^2 = \frac{r^2}{\para{\abs{c}^2 - r^2}^2}
            \end{align*}
            is a Euclidean circle.
        \end{itemize}
        \item The case where \(C\) is a line is similar. 
    \end{itemize}
\end{proof}

We now relate the cross ratio defined in \autoref{def:cross-ratio} to geometry.

\begin{corollary}[Condition for Points in \(\CC\compactify\) to lie on a Circle]
    \(4\) points \(z_1, z_2, z_3, z_4 \in \CC\compactify\) lie on a circle if and only if their cross ratio
    \[
        \brac{z_1, z_2, z_3, z_4} \in \RR\compactify = \RR \cup \brce{\infty}.
    \]
\end{corollary}

\begin{proof}
    From \autoref{def:cross-ratio}, let \(\alpha \in \mobius\) such that
    \begin{align*}
        \alpha \para{z_1} &= 0, \\
        \alpha \para{z_2} &= 1, \\
        \alpha \para{z_3} &= \infty, \\
        \alpha \para{z_4} &= \brac{z_1, z_2, z_3, z_4}.
    \end{align*}

    We prove the two directions separately.
    \begin{itemize}
        \item \(\implies\). If \(z_1, z_2, z_3, z_4\) lie on a circle \(C\), then by \autoref{thm:mobius-preserve-circles}, \(0, 1, \infty, \alpha\para{z_4}\) lie on a circle \(\alpha\para{C}\).
        
        But the only circle containing \(0, 1, \infty\) is the line \(\RR\compactify\), and hence \(\alpha\para{z_4} \in \RR\compactify\) as required.

        \item \(\impliedby\). If \(\brac{z_1, z_2, z_3, z_4} \in \RR\compactify\), then \(\alpha\para{z_1}, \alpha\para{z_2}, \alpha\para{z_3}, \alpha\para{z_4} \in \RR\compactify\) is a circle, so
        \[
            z_1, z_2, z_3, z_4 \in \alpha\inverse\para{\RR\compactify}
        \]
        which is a circle as well, by \autoref{thm:mobius-preserve-circles}.
    \end{itemize}
\end{proof}